\documentclass{article}
\usepackage{amsmath}
\usepackage{amssymb}
\setlength{\parindent}{0pt}

\begin{document}

% ======================== SOAL 1 ========================
\textbf{1.} $\displaystyle \int \left(12x^2 - 4x + 1\right) dx = \ldots$

\medskip
\textbf{Jawab:}

Pakai aturan pangkat: $\displaystyle \int x^n \, dx = \frac{x^{n+1}}{n+1} + C$

\begin{align*}
  \int 12x^2 \, dx &= 12 \cdot \frac{x^{3}}{3} = 4x^3 \\[6pt]
  \int -4x \, dx   &= -4 \cdot \frac{x^{2}}{2} = -2x^2 \\[6pt]
  \int 1 \, dx     &= x
\end{align*}

\begin{align*}
  \int \left(12x^2 - 4x + 1\right) dx &= 4x^3 - 2x^2 + x + C
\end{align*}

Cek (turunkan balik):
\begin{align*}
  \frac{d}{dx}\left(4x^3 - 2x^2 + x + C\right) &= 12x^2 - 4x + 1 \quad \checkmark
\end{align*}

\boxed{4x^3 - 2x^2 + x + C}

\bigskip\hrule\bigskip

% ======================== SOAL 2 ========================
\textbf{2.} $\displaystyle \lim_{x \to 0} \frac{e^x - 1 - x}{x^2} = \ldots$

\medskip
\textbf{Jawab:}

Substitusi $x = 0$:
\begin{align*}
  \frac{e^0 - 1 - 0}{0^2} &= \frac{0}{0} \quad \text{(tak tentu} \rightarrow \text{L'H\^opital)}
\end{align*}

\begin{align*}
  \lim_{x \to 0} \frac{e^x - 1 - x}{x^2}
    &= \lim_{x \to 0} \frac{e^x - 1}{2x}
       && \text{(masih } \tfrac{0}{0} \text{, L'H\^opital lagi)} \\[6pt]
    &= \lim_{x \to 0} \frac{e^x}{2} \\[6pt]
    &= \frac{e^0}{2} \\[6pt]
    &= \frac{1}{2}
\end{align*}

\boxed{\dfrac{1}{2}}

\bigskip\hrule\bigskip

% ======================== SOAL 3 ========================
\textbf{3.} $\displaystyle \lim_{x \to 0^+} x^x = \ldots$

\medskip
\textbf{Jawab:}

Bentuknya $0^0$ (tak tentu).

Misal $y = x^x$, maka
\begin{align*}
  \ln y &= \ln x^x = x \ln x \quad \text{(bentuk } 0 \cdot (-\infty)\text{)}
\end{align*}

Ubah jadi pecahan:
\begin{align*}
  \lim_{x \to 0^+} \ln y
    &= \lim_{x \to 0^+} \frac{\ln x}{\frac{1}{x}}
       && \text{(bentuk } \tfrac{-\infty}{\infty} \rightarrow \text{L'H\^opital)} \\[6pt]
    &= \lim_{x \to 0^+} \frac{\frac{1}{x}}{-\frac{1}{x^2}} \\[6pt]
    &= \lim_{x \to 0^+} \frac{1}{x} \cdot \left(-x^2\right) \\[6pt]
    &= \lim_{x \to 0^+} (-x) \\[6pt]
    &= 0
\end{align*}

Karena $\ln y \to 0$:
\begin{align*}
  y &\to e^0 = 1
\end{align*}

\boxed{1}

\bigskip\hrule\bigskip

% ======================== SOAL 4 ========================
\textbf{4.} Tentukan konvergen/divergen: $\displaystyle \int_1^{\infty} \frac{\ln x}{x^2}\, dx$

\medskip
\textbf{Jawab:}

Batas atas $\infty$ $\rightarrow$ integral tak wajar, ubah ke limit:
\begin{align*}
  \int_1^{\infty} \frac{\ln x}{x^2}\, dx &= \lim_{b \to \infty} \int_1^{b} \frac{\ln x}{x^2}\, dx
\end{align*}

Integral parsial $\displaystyle \int u \, dv = uv - \int v \, du$

\begin{align*}
  u  &= \ln x    & dv &= x^{-2}\, dx \\
  du &= \frac{1}{x}\, dx & v &= -\frac{1}{x}
\end{align*}

\begin{align*}
  \int \frac{\ln x}{x^2}\, dx
    &= \ln x \cdot \left(-\frac{1}{x}\right) - \int \left(-\frac{1}{x}\right)\frac{1}{x}\, dx \\[6pt]
    &= -\frac{\ln x}{x} + \int \frac{1}{x^2}\, dx \\[6pt]
    &= -\frac{\ln x}{x} - \frac{1}{x}
\end{align*}

Masukkan batas:
\begin{align*}
  \int_1^{b} \frac{\ln x}{x^2}\, dx
    &= \left[-\frac{\ln x}{x} - \frac{1}{x}\right]_1^b \\[6pt]
    &= \left(-\frac{\ln b}{b} - \frac{1}{b}\right) - \left(-\frac{\ln 1}{1} - 1\right) \\[6pt]
    &= -\frac{\ln b}{b} - \frac{1}{b} + 1 \qquad (\ln 1 = 0)
\end{align*}

Cari limitnya, $\dfrac{\ln b}{b}$ bentuk $\dfrac{\infty}{\infty}$ $\rightarrow$ L'H\^opital:
\begin{align*}
  \lim_{b \to \infty} \frac{\ln b}{b} &= \lim_{b \to \infty} \frac{\frac{1}{b}}{1} = 0
\end{align*}

\begin{align*}
  \lim_{b \to \infty} \left(-\frac{\ln b}{b} - \frac{1}{b} + 1\right) &= -0 - 0 + 1 \\
  &= 1
\end{align*}

Limitnya ada (hingga), jadi integralnya \textbf{konvergen}.

\boxed{\displaystyle \int_1^{\infty} \frac{\ln x}{x^2}\, dx = 1 \;\text{(konvergen)}}

\end{document}